\section{研究结论}

本文提出FinPlanRAG，把金融混合查询的回答形式选择表述为查询契约路由，并以义务检索和AnchorEvolve分别处理确定信息与前瞻性断言。路由器在检索与生成之前判断事实是否已在截止时点存在；义务检索围绕类型化义务形成闭合答案；AnchorEvolve以L0稳定关系、L1行业模式和L2实例锚点维护可更新、可追溯的前瞻性证据状态。校准读出在此基础上生成点答案、可信集或弃权，用于控制最终决策风险。

第一项贡献是查询契约路由。现有问答、断言验证和选择性预测分别固定输出接口，本文把“事实是否已在截止时点存在”置于检索与生成之前，用一次路由决定闭合答案或可校准判断。路由实验支持这一贡献：DeepSeek Flash和GLM 5.3 Flash分别达到0.995和0.990的路由准确率，混合效用接近oracle，并显著高于Always-RAG和Always-Anchor；Qwen3-8B的结果则说明路由能力依赖基座。契约路由的价值不只是提高一个分类指标，而是把混合负载的分析单位从数据集任务改为查询契约，使路由错误和分支错误可以分开诊断。

第二项贡献是义务检索与AnchorEvolve的两条证据分支。义务检索把确定问题绑定为类型化义务，在时点语料上按覆盖取证，由阅读器返回带证据编号的闭合答案；AnchorEvolve把前瞻性断言映射到有限结果空间，用三层锚点、五类资格判定和时点信念更新保留证据的作用域、可靠性和时效。DocFinQA和FinDVer实验支持义务检索的效率主张；收入指引实验和锚点消融支持AnchorEvolve的语义主张。特别是去除全部锚点后LLM后验宏F1下降，说明锚点不是装饰，而是改变了模型判断的中间状态。

第三项贡献是语义状态与校准读出的分工，以及统一的多基准评估。本文在三个数据集、三套基座和六个具名基线上同时检验路由、阅读、锚点推断、后验融合、校准选择和外部迁移，并把分类、校准、证据成本、决策损失和路由效用放在同一口径下报告。需要明确，纯RF是当前最强的数值决策头，融合后验在部分基座上取得更高普通准确率，但决策损失仍高于纯RF；因此本文不主张LLM锚点后验单独胜过强RF。本文主张的是：LLM在路由、阅读和语义状态构造上不可替代，校准读出负责把这些状态转换为低风险决策，二者应当分开评估。

本文的原创性贡献可以概括为三点。第一，把输出契约从固定接口变成查询层的路由变量，使确定问答与前瞻性裁决在同一时点约束下被统一处理。第二，把前瞻性证据从单一标签变成可审计的状态，使证据来源、作用域、资格理由和信念变化可以被逐条检查。第三，把语义判断与风险校准分开，使LLM的语义能力、统计模型的风险控制能力和路由器的契约识别能力可以分别检验、替换和组合。这三点不依赖特定基座或特定行业，构成了本文相对于现有金融RAG和断言验证工作的主要区别。

从金融实践看，框架把答案、证据和风险放在同一条审计链上。分析师对确定问题获得带证据编号的闭合答案，对前瞻性断言获得点判断、可信集或弃权，并可以查看锚点来源和资格理由。机构可以按自身风险偏好替换损失参数和校准后端，而不改变路由与证据状态的结构。生产里程碑、产能投放、监管审批和合同履约等任务与收入指引具有相似的时点结构，本文结果为这些场景提供可复用的接口。

本文的结论有明确边界。收入指引实验采用三桶结果空间，通用的时机与阶段状态仍需在更多前瞻性事件上评估；混合负载由DocFinQA与收入指引样本构成，真实用户查询的语义边界可能更复杂；标签来自XBRL核验而非人工金标准。后续工作将围绕更多前瞻性事件、更稳定的边界路由、更窄的可信集和更严格的人工复核展开。由于路由、锚点抽取和输出读出彼此分离，这些改进可以独立评估，而不需要改变FinPlanRAG的整体接口。
